% !TEX root = ../main.tex
\chapter{Introduction}

International remittances are one of the main private insurance channels through which local economic shocks can cross national borders \parencite{yang_choi_2007, bettin_2023, dorantes_pozo_2006}. In 2024, Mexico received more than US\$67 billion in personal remittances, making it the second largest receipient country in the world after India \parencite{banxico_remittances_ca11, worldbank_remittances}. The U.S.--Mexico corridor is especially important because 97\% of the migrant population in Mexico resides in the United States \parencite{LiNg2024}. This paper asks whether a negative local shock in the United States is transmitted to Mexican municipalities through the lens of remittances. Do migrants smooth transfers after being hit by a shock, or do origin households bear part of the adjustment through lower remittance inflows?
\par
The answer matters because remittances are a significant part of household resources, with migrant households in the U.S. of Mexican origin sending on average around US\$12,000 per year \parencite{isidro_martinez_2025}. While our data does not directly observe Mexican household income, prior work shows that remittances in Mexico affect consumption, poverty, inequality, food security, schooling, and child labor decisions \parencite{adamsetal2005,MORARIVERA2021105349, ALCARAZ2012156}. We therefore interpret changes in municipal remittance inflows as reduced-form evidence on a channel with direct relevance for household income and consumption, while not claiming to estimate those household outcomes themselves.
\par
Our paper contributes to two distinct strands of the migration and remittance literature. First, we contribute to the literature on international shock transmission, for which the primary empirical difficulty liest in the fact that major economic crises typically affect both sides of a migration corridor. For instance, while \textcite{caballero_2023} provide important evidence that U.S. labor market shocks during the Great Recession reduced remittances to Mexico, the global nature of that crisis affected both countries simultaneously. Furthermore, their study focuses on the extensive margin of remittance receipt. We overcome these challenges by exploiting Hurricane Ian, the costliest hurricane in Florida's history \parencite{nhc_ian_2026}, which struck solely the sender side of the corridor. Because natural disasters severely restrict local labor demand and household resources \parencite{natural_disasters_US}, Hurricane Ian provides a clean, exogenous contraction to migrants' remitting capacity. This allows us to quantify the intensive margin of bilateral remittance adjustments, while also studying the effect of a shock with much more limited aggregate consequences than a global recession.
\par
Second, we contribute to the literature studying remittances as a form of informal household insurance \parencite{yang_choi_2007, bettin_2023, dorantes_pozo_2006}. The vast majority of existing evidence focuses on shocks in the receiving country, demonstrating how migrants step up transfers to insure their families against local disasters at origin. We invert this framework by examining a sender-side shock, where the economic disruption occurs in the destination labor market while origin municipalities remain untouched.
\par
Chapter \ref{chap:structural_model} introduces a structural remittance model to formalize the impact of the shock on remittance flows. The model follows the insight in \textcite{albert_monras_2022} that migrants allocate resources across where they live and where their origin households are located. In our setting, this idea is applied to remittance choices rather than to location choices. Migrants choose between local consumption in the United States and resources available to the origin household in Mexico. They differ in income and in the strength of their altruistic or obligation motive, and remitting involves both a fixed cost and a variable cost. In the spirit of models with heterogeneous agents, fixed costs, and variable costs \parencite{melitz2003}, the fixed cost generates an extensive margin of remitting, while the variable cost governs the intensity of transfers. We then augment the model to account for the fact that origin-household resources include local non-remittance income plus remittances received from migrants in other U.S. states. This extension allows us to interpret both the direct exposure effect and possible spillover effects \parencite{hudgens2008, aronow2017, butts2023}.
\par
Chapter \ref{chap:data_empirics} connects the model to the data by constructing the inflow of remittances flows from U.S. states to Mexican municipalities. Because such flows are not directly observed, we combine two sources from Banco de M{\'e}xico with MCAS administrative records. Banco de M{\'e}xico reports remittance inflows by U.S. state of origin and by Mexican municipality of receipt \parencite{banxico_ce168, banxico_ce166}. MCAS records identify the Mexican municipality of origin and U.S. state of residence of Mexican consular-card holders, providing a proxy for the geography of migrant networks. We use these records to construct state-to-municipality migration shares and allocate observed state-level remittance outflows across Mexican municipalities. We focus our analysis on the Mexican municipalities receiving remittances from Florida.

Two validation exercises support this construction. First, the migration network data are highly persistent: the distribution of Mexican migrants across U.S. states and Mexican municipalities has year-to-year correlations above 0.90 from 2012 onward. Second, when we use these migration shares to allocate state-level remittance outflows to Mexican municipalities, the resulting municipal aggregates correlate between 0.75 and 0.85 with official municipal remittance receipts. These facts suggest that the estimated corridors capture meaningful variation in the geography of remittances and that the assumptions underlying our construction are reasonable.

The third part of the paper studies the shock. We exploit variation in Mexican municipalities' pre-existing exposure to Florida to estimate whether Ian affected municipal remittance inflows differently depending on their exposure to the shocked state. In the period immediately following Ian, a 1\% increase in a municipality's pre-existing exposure to Florida is associated with a 0.16\% decrease in remittances received. This result is consistent with a sender-side income shock being transmitted through the remittance corridor to origin municipalities.

Finally, we examine migration-network spillovers. Among municipalities jointly exposed to Florida and other major sender states, remittances originating from California remain largely unchanged following the shock. By contrast, remittances from Texas, the second-largest sender state, decline in municipalities that are jointly exposed to both migration corridors in the post-Ian period. We explore this relationship further by centering the data on different quantiles of Texas exposure to Mexican municipalities. The diagnostic plots show for the 10\% most exposed Florida municipalities to Texas, (insert write-up here), whereas for the least 10\% exposed show (describe what we see). This pattern is consistent with both compensatory transfers from migrants in other states and out-migration from Florida-exposed counties.

